\section{The \texttt{Pol\_tabled\_field} McStas Component}
Magnetic field component.

\subsection*{Identification}
\begin{itemize}
  \item \textbf{Site:} 
  \item \textbf{Author:} Erik B Knudsen, Peter Christiansen and Peter Willendrup
  \item \textbf{Origin:} RISOE
  \item \textbf{Date:} July 2011
\end{itemize}

\subsection*{Description}
\begin{lstlisting}
Region with a tabled magnetic field read from file.
The magnetic field is read from a text file where it is
specified as a point cloud with N rows of 6 columns:
x y z Bx By Bz
the B field map is resampled with Stepx*Stepy*Stepz points.
Use Stepx=Stepy=Stepz=0 to skip resampling and use the table as is.
The regions itself may be either a 3D rectangular block, a cylinder with axis along y,
or spherical. Interpolation is done between data-points.
\end{lstlisting}

\subsection*{Input parameters}
Parameters in \textbf{boldface} are required; the others are optional.

\begin{longtable}{p{0.22\textwidth}p{0.12\textwidth}p{0.46\textwidth}p{0.14\textwidth}}
\toprule
\textbf{Name} & \textbf{Unit} & \textbf{Description} & \textbf{Default} \\
\midrule
\endhead
xwidth & m & Width of opening. & 0 \\
yheight & m & Height of opening. & 0 \\
zdepth & m & Length of field. & 0 \\
radius & m & Radius of field if it is cylindrical or spherical. & 0 \\
filename & str & File where the magnetic field is tabulated. & "bfield.dat" \\
geometry & str & Name of an Object File Format (OFF) or PLY file for complex field-geometry. & NULL \\
interpol\_method & str & Choice of interpolation method "kdtree" (default on CPU) / "regular" (default on GPU) & "default" \\
\bottomrule
\end{longtable}

\subsection*{Links}
\begin{itemize}
  \item \href{run:/home/willend/willend-McCode/mcstas-comps/optics/Pol_tabled_field.comp}{Source code} for \texttt{Pol\_tabled\_field.comp}.
\end{itemize}
\IfFileExists{Pol_tabled_field_static.tex}{\input{Pol_tabled_field_static.tex}}{}